\subsection{Encumbrances}

DEFINITION 1. --- Let T be a set. One calls encumbrance on T any mapping p of \(\mathcal{F}_{+}(\mathrm{T})\) into \(\overline{\mathbf{R}}_{+}\) that has the following properties:

\begin{enumerate}
\def\labelenumi{\alph{enumi})}
\item
  If \(f\) and \(g\) are two elements of \(\mathcal{F}_+\) such that \(f \leqslant g\) , then \(p(f) \leqslant p(g)\) .
\item
  If \(f\) is an element of \(\mathcal{F}_+\) , and \(t\) is a number \(\geqslant 0\) , then \(p(tf) = tp(f)\) .
\item
  If \(f\) and \(g\) are two elements of \(\mathcal{F}_+\) , then \(p(f + g) \leqslant p(f) + p(g)\) .
\item
  If \((f_n)\) is an increasing sequence of elements of \(\mathcal{F}_+\) , and if \(f = \lim_{n \to \infty} f_n\) , then \(p(f) = \lim_{n \to \infty} p(f_n)\) .
\end{enumerate}

If A is a subset of T, we write \(p(A)\) instead of \(p(\varphi_A)\) .

The condition \(b)\) implies that \(p(0) = 0\) . On the other hand, let \((f_n)\) be a sequence of elements of \(\mathcal{F}_+\) ; the conditions \(c)\) and \(d)\) imply the inequality

\[
p \left(\sum_ {n} f _ {n}\right) \leqslant \sum_ {n} p (f _ {n})
\]

(the inequality of countable convexity).

For example, let T be a locally compact space, \(\mu\) a positive measure on T; then \(\mu^{*}\) and \(\mu^{\bullet}\) are encumbrances on T. This follows from Props. 10, 11, 12 and Th. 3 of Ch. IV, §1, No.~3 for \(\mu^{*}\) , and from Prop. 1 of Ch. V, §1, No.~1 for \(\mu^{\bullet}\) .

PROPOSITION 1. --- Let \((p_{\alpha})_{\alpha \in A}\) be a family of encumbrances on T. The sum and upper envelope of the family \((p_{\alpha})\) (in \(\mathcal{F}_{+}(\mathcal{F}_{+}(T))\) ) are then encumbrances.

The sum of a finite family of encumbrances obviously being an encumbrance, it suffices to treat the case of the upper envelope. The properties \(a), b), c)\) of Definition 1 being obviously satisfied, it remains to establish \(d)\) . Set \(p = \sup p_{\alpha}\) ; then, with the notations of Definition 1 \(d)\) ,

\[
p (f) = \sup _ {\alpha} p _ {\alpha} (f) = \sup _ {\alpha} \sup _ {n} p _ {\alpha} (f _ {n}) = \sup _ {n} \sup _ {\alpha} p _ {\alpha} (f _ {n}) = \sup _ {n} p (f _ {n}).
\]

DEFINITION 2. --- Let p be an encumbrance on a set T. One says that p is bounded if \(p(T) < +\infty\) . If T is a topological space, p is said to be locally bounded provided that every \(x \in T\) admits a neighborhood V such that \(p(V) < +\infty\) .

It then follows from the properties \(a\) and \(c\) of Def. 1 that \(p(\mathbf{K}) < +\infty\) for every compact subset \(\mathbf{K}\) of \(\mathbf{T}\) . In particular, if \(\mathbf{T}\) is compact, then every locally bounded encumbrance on \(\mathbf{T}\) is bounded.

Let \(p\) be an encumbrance on a set \(T\) , and \(A\) a subset of \(T\) . For every function \(f \in \mathcal{F}_{+}(A)\) , let \(f^0\) be the extension by 0 of \(f\) to \(T\) ; the mapping \(f \mapsto p(f^0)\) on \(\mathcal{F}_{+}(A)\) is then an encumbrance, called the encumbrance induced by \(p\) on \(A\) , and is denoted \(p|_{A}\) or \(p_A\) .

Let T and U be two sets, \(\pi\) a mapping of T into U, and p an encumbrance on T. The encumbrance \(\pi(p)\) on U, whose value for \(f \in \mathcal{F}_{+}(U)\) is given by

\[
\bigl (\pi (p) \bigr) (f) = p (f \circ \pi),
\]

is called the image encumbrance of \(p\) under \(\pi\) .

Let \(p\) be an encumbrance on a set \(T\) ; \(p\) is said to be concentrated on a subset \(A\) of \(T\) if \(p(T - A) = 0\) .

Lemma 1. --- If the encumbrance \(p\) is concentrated on \(\mathbf{A} \subset \mathbf{T}\) , then \(p(f) = p(f\varphi_{\mathbf{A}})\) for every \(f \in \mathcal{F}_{+}(\mathbf{T})\) .

For, set \(\mathbf{T} - \mathbf{A} = \mathbf{B}\) , so that \(p(\varphi_{\mathrm{B}}) = 0\) ; then

\[
f \varphi_ {\mathrm{B}} \leqslant (+ \infty) \cdot \varphi_ {\mathrm{B}} = \sup _ {n \in \mathbf {N}} n \varphi_ {\mathrm{B}},
\]

therefore \(p(f\varphi_{\mathrm{B}}) = 0\) by properties \(a), b), d)\) of Def. 1. It follows from \(c)\) that \(p(f) \leqslant p(f\varphi_{\mathrm{A}}) + p(f\varphi_{\mathrm{B}}) = p(f\varphi_{\mathrm{A}})\) , and finally \(p(f) = p(f\varphi_{\mathrm{A}})\) by \(a)\) .

